
\section{Attached File}

We provide four new examples that demonstrate the expressivity
of \PAT{}s (Reviewer E):
\begin{enumerate}
\item "exactly one put event"

\item sequential consistency (i.e., a behavior involving multiple
  logical threads)

\item deadlock freedom

\item network behaviors
\end{enumerate}
We also include all the \PAT{}s used in our experimental evaluation,
starting with the two requested by the reviewers (2PC and
RingLeaderElection).

\subsection{Additional Examples}

\paragraph{Exactly one put event} We can add a request ID to each
message via a \(\I{rid}\) field and then constrain it to be unique:
{\small\begin{align*}
    i{:}\Code{rId} \sarr \rg{H}{\lastA \msg{putRsp}{rid = i}}{F \land \neg \finalA (\msg{putRsp}{rid = i} \land \nextA \finalA \msg{putRsp}{rid = i})}
\end{align*}}

\paragraph{Sequential consistency} We demonstrate how to encode
a sequential consistency property with a simple example where
multiple nodes (processes) perform read and write operations on a
database. Each node communicates with the database using
\(\eff{readReq/Rsp}\) and \(\eff{writeReq/Rsp}\) events, which
represent the local view of each process. For each of these
operations, the controller also sends duplicate auxiliary events
\(\eff{gReadReq/Rsp}\) and \(\eff{gWriteReq/Rsp}\). Then sequential
consistency can be stated as: %
\begin{enumerate}
\item The global order matches the local view:
  {\small\begin{align*}
           i{:}\Code{requestId} \garr &(\msg{writeReq}{\I{rid} = i} \implies \nextA \msg{gWriteReq}{\I{rid} = i}) \land
           \\&(\msg{writeRsp}{\I{rid} = i} \implies \nextA \msg{gWriteRsp}{\I{rid} = i}) \land ...
         \end{align*}}\noindent
       which requires that local operations are immediately matched by an
       auxillary operation.

     \item The execution remains the same as if all operations were
       performed in the global order, which can be captured in the
       types of $\eff{readReq}$ and $\eff{readRsq}$:
       {\small\begin{align*}
        \eff{readReq}&: i{:}\Code{rId} \garr \Code{x}\hasnty{int} \sarr
        \\&\rg{\finalA (\msg{gWriteReq}{v = \Code{x}} \land \nextA \neg \finalA \eff{gWriteReq})}{\lastA \msg{readReq}{\I{id} = i}}{\msg{readRsp}{v = \Code{x} \land \I{id} = i}}
        \\\eff{readRsq}&: i{:}\Code{rId} \sarr \Code{x}\hasnty{int} \sarr
        \\&\rg{\finalA (\msg{writeReq}{v = \Code{x}} \land \nextA ((\neg \eff{writeReq}) \untilA \msg{readRsp}{v = \Code{x} \land \I{id} = i })) }{\lastA \msg{readRsp}{\I{id} = i}}{\allTL}
    \end{align*}}\noindent
  where the type of $\eff{readReq}$ requires the value read to be
  last written value according to the controller's view
  ($\eff{gWriteReq}$), and the \PAT{} of $\eff{readRsp}$ requires it
  to be the last value written under the actor's view
  ($\eff{writeReq}$).

\end{enumerate}

\paragraph{Deadlock Freedom}

We demonstrate how to encode deadlock freedom using a set of nodes,
each of which requires holding some resources before it can make
progress. A node \( n \) will use \(\eff{aquire}(n, r)\) and
\(\eff{release}(n, r)\) events to request and release a resource
\( r \). A controller will use the \(\eff{stuck}(n, r)\) to record
that node \( n \) is waiting to claim a resource \( r \) (i.e., it
fails to \(\eff{aquire}\) the resource \( r \)). The controller uses
the \(\eff{invoke}(n)\) message to notify node \( n \) that the
corresponding resource is now available. Deadlock occurs when all
nodes \( m \) are stuck on some resource:
{\small\(\finalA (\msg{stuck}{\I{n} = m} \land
  \neg\finalA\msg{invoke}{\I{n} = m})\)}.

\paragraph{Network Behavior}
Suppose an actor sends message \(\eff{E_1}\) followed by
\(\eff{E_2}\). Then, it can stipulate in-order delivery of these
messages via the \PAT{} {\small\(\rg{H}{A}{\finalA (\eff{E_1} \land \nextA \finalA
    \eff{E_2})}\)}, while an out-of-order delivery is expressed as
{\small\(\rg{H}{A}{(\finalA \eff{E_1}) \land (\finalA \eff{E_2})}\)}.
Node failures can be represented as a special event, \(\eff{crash}\),
sent by the controller; the following \PAT{} states that an actor can
only process messages if it has not received a \(\eff{crash}\) event:
{\small\(\rg{H \land \neg (\finalA \eff{crash})}{A}{F}\)}.

\subsection{Benchmark Explanation}

\paragraph{\textsc{RingLeaderElection}} The benchmark involves three
events:
\begin{enumerate}

\item The message \(\eff{nominate}(\I{s},\I{d},\I{l})\) indicates that
  node \(\I{s}\) sends a message to node \(\I{d}\) to nominate node
  \(\I{l}\) as the leader;

\item The message \(\eff{won}(\I{l})\) announces
  that node \(\I{l}\) is the leader;

\item The message \(\eff{wakeup}(\I{n})\) enables node \(\I{n}\) to
  elect a leader.
\end{enumerate}
  The unique leader property requires that no two
  different nodes are announced as the leader. The \PAT{}s used in
  this benchmark are:
%
  {\small
    \begin{align*}
    \eff{nominate}&: \Code{n}\hasoty{tNode}{\top} \sarr \Code{m}\hasoty{tNode}{\top}
  \sarr \Code{ld}\hasoty{tNode}{\top} \sarr
  \\&\quad \rg{\globalA\evparenth{\top}}{\lastA\msg{nominate}{
      \I{s} = \Code{n} \land \I{d} = \Code{m} \land \I{l} =
      \Code{ld} \land \Code{m = ld} }}{\finalA\msg{won}{\I{l} =
      \Code{m}}} \interty
  \\&\quad \rg{\globalA\evparenth{\top}}{\lastA\msg{nominate}{\I{s} = \Code{n} \land \I{d} = \Code{m} \land \I{l} = \Code{ld} \land \Code{m \not= ld} }}{\finalA\msg{nominate}{\I{s} = \Code{m} \land \I{d} = \I{next}(\Code{m}) \land \I{l} = \Code{ld}}}
  \\\eff{won}&: \Code{n}\hasoty{tNode}{\top} \sarr \rg{\allTL}{\msg{won}{\I{l} = n}}{\allTL}
  \\\eff{wakeup}&: \Code{m}{:}\Code{tNode}\sarr \rg{\allTL}{\msg{wakeup}{\I{n} = \Code{m}}}{\finalA \msg{nominate}{\I{s} = \Code{m} \land \I{d} = \I{next}(\Code{m}) \land \I{l} = \Code{m}}}
\end{align*}}\noindent

The predicate \(\I{next}(n_1, n_2)\) indicates that node \(n_2\) is
the next node of \(n_1\) in the ring. The two intersected \PAT{}s of
\(\eff{nominate}\) describe two situations in which a node can receive
a nomination message:

\begin{enumerate}
    \item If the nominated leader has an equal ID, announce itself as the leader.
    \item If the nominated leader has a different ID, nominate this leader.
\end{enumerate}

The protocol is incorrect with respect to the unique leader property
because a node should only confirm a previous nomination if the
nominated leader has a greater ID, rather than a different ID.

\paragraph{\textsc{Simplified2PC}}

In this simplified version of the two-phase commit protocol, we assume
transactions consist of a single update operation. The controller
communicates with the coordinator via \(\eff{read/write}\) messages,
and the coordinator communicates with the underlying database via
\(\eff{put/get}\) messages. Depending on the response from the
database, the coordinator can \(\eff{commit/abort}\) a transaction.
The \PAT{}s of \(\eff{readRsp/writeRsp/commit/abort}\) are not
particularly interesting, as handling these messages does not result
in new messages being sent; thus, we omit them. The \PAT{}s of rest
operators are defined as the following: %

{\small\begin{align*}
  \eff{readReq}&: \rg{\allTL}{\lastA\msg{readReq}{\top}}{\finalA\msg{getReq}{\top}}
  \\\eff{getReq}&: x{:}\Code{tVal}\garr
  \\&\rg{\finalA(\msg{putReq}{v = x} \land \nextA(\finalA \msg{commit}{\top} \land \finalA\msg{putReq}{\top}))}{\lastA\msg{getReq}{\top}}{\finalA\msg{getRsp}{v = x}} \interty
  \\&\rg{\neg \finalA \msg{commit}{\top}}{\lastA\msg{getReq}{\top}}{\finalA\msg{getRsp}{v = -1}}
  \\\eff{writeReq}&: x{:}\Code{tVal}\sarr \rg{\allTL}{\msg{writeReq}{v= x}}{\finalA\msg{putReq}{v = x}}
  \\\eff{putReq}&: x{:}\Code{tVal}\sarr \rg{\allTL}{\msg{putReq}{v= x}}{\finalA\msg{putRsp}{v = x}}
  \\\eff{putRsp}&: x{:}\Code{tVal}\sarr s{:}\Code{status}\sarr 
  \\&\rg{\allTL}{\msg{putRsp}{v= x \land \I{stat} = s \land s}}{\finalA\msg{writeRsp}{v = x \land \I{stat}} \land \finalA\msg{commit}{\top}} \interty
  \\&\rg{\allTL}{\msg{putRsp}{v= x \land \I{stat} = s \land \neg s}}{\finalA\msg{writeRsp}{v = x \land \neg \I{stat}} \land \finalA\msg{abort}{\top}}
\end{align*}}\noindent

Note that the database handles the \(\eff{getReq}\) message by
returning the default value (\(-1\)) when no commit has occurred yet;
otherwise, it returns the last committed value. The coordinator
handles the \(\eff{getReq}\) message by sending a write response to
the controller and \emph{simultaneously} sending a commit/abort to the
database, which leads to a violation of the RYW property.

\paragraph{\textsc{Database}} This benchmark was used as the
motivating example in the paper.

\paragraph{\textsc{Firewall}} A set of internal and external nodes
communicate through a firewall. The firewall should block messages
from an external node unless that node has previously received a
message from an internal node. In practice, the firewall maintains a
whitelist of external nodes that are permitted to communicate with
internal nodes. {\small\begin{align*}
  \eff{internalReq}&: \Code{n}\hasoty{tNode}{\I{internal}(\vnu)} \sarr \Code{m}\hasoty{tNode}{\I{external}(\vnu)} \sarr
  \\&\quad\rg{\allTL}{\lastA\msg{internalReq}{\I{s} = n \land \I{d} = m}}{\finalA\msg{forwardReq}{\I{s} = n \land \I{d} = m}}
  \\\eff{forwardReq}&: \Code{n}\hasoty{tNode}{\I{internal}(\vnu)} \sarr \Code{m}\hasoty{tNode}{\I{external}(\vnu)} \sarr
  \\&\quad\rg{\allTL}{\lastA\msg{forwardReq}{\I{s} = n \land \I{d} = m}}{\finalA\msg{externalReq}{\I{s} = m \land \I{d} = n}}
  \\\eff{externalReq}&: \Code{m}\hasoty{tNode}{\I{external}(\vnu)} \sarr \Code{n}\hasoty{tNode}{\I{internal}(\vnu)} \sarr
  \\&\quad
  \effLR{\finalA \msg{internalReq}{\I{d} = m}\land \neg\nextA\finalA\msg{internalReq}{\top}}
  \effLR{\lastA\msg{externalReq}{\I{s} = n \land \I{d} = m}}\\&\qquad\effLR{\finalA\msg{externalRsp}{\I{s} = n \land \I{d} = m}}\interty
  \\&\quad\rg{\neg(\finalA \msg{internalReq}{\I{d} = m}\land \neg\nextA\finalA\msg{internalReq}{\top})}{\lastA\msg{externalReq}{\I{s} = n \land \I{d} = m}}{\allTL}
  \\\eff{externalRsp}&: \Code{m}\hasoty{tNode}{\I{external}(\vnu)} \sarr \Code{n}\hasoty{tNode}{\I{internal}(\vnu)} \sarr
  \rg{\allTL}{\lastA\msg{externalRsp}{\I{s} = n \land \I{d} = m}}{\allTL}
\end{align*}}\noindent

Note that this whitelist only stores the most recent external node (as
specified in the first \PAT of \(\eff{externalReq}\)), which violates
the (informal) liveness property that all messages should eventually
receive a response.

\paragraph{\textsc{EspressoMachine}} The user interacts with a coffee machine through its control panel, where the panel must correctly interpret user inputs and handler errors from coffee machine.
{\small\begin{align*}
  \eff{eCoffeeMachineUser}&: \rg{\allTL}{\lastA\msg{eCoffeeMachineUser}{\top}}{\finalA\msg{eWarmUpReq}{\top}}
  \\\eff{eEspressoButtonPressed}&: \effLR{\finalA \msg{eCoffeeMakerReady}{\top} \land \neg\nextA\finalA \msg{eCoffeeMakerError}{\top}}
  \\&\quad \effLR{\lastA\msg{eEspressoButtonPressed}{\top}}\effLR{\finalA\msg{eGrindBeansReq}{\top}}
  \\\eff{eCoffeeMakerError}&: \rg{\allTL}{\lastA\msg{eCoffeeMakerError}{\top}}{\allTL}
  \\\eff{eCoffeeMakerReady}&: \rg{\allTL}{\lastA\msg{eCoffeeMakerReady}{\top}}{\allTL}
  \\\eff{eWarmUpReq}&: \rg{\allTL}{\lastA\msg{eWarmUpReq}{\top}}{\finalA\msg{eWarmUpCompleted}{\top}}
  \\\eff{eGrindBeansReq}&: \rg{\allTL}{\lastA\msg{eGrindBeansReq}{\top}}{\finalA\msg{eNoBeansError}{\top}} \interty
  \\&\quad \rg{\allTL}{\lastA\msg{eGrindBeansReq}{\top}}{\finalA\msg{eGrindBeansCompleted}{\top}}
  \\\eff{eStartEspressoReq}&: \rg{\allTL}{\lastA\msg{eStartEspressoReq}{\top}}{\finalA\msg{eNoWaterError}{\top}} \interty
  \\&\quad \rg{\allTL}{\lastA\msg{eStartEspressoReq}{\top}}{\finalA\msg{eEspressoCompleted}{\top}}
  \\\eff{eGrindBeansCompleted}&: \rg{\allTL}{\lastA\msg{eGrindBeansCompleted}{\top}}{\finalA\msg{eStartEspressoReq}{\top}}
  \\\eff{eEspressoCompleted}&: \rg{\allTL}{\lastA\msg{eEspressoCompleted}{\top}}{\finalA\msg{eCoffeeMakerCompleted}{\top}}
  \\\eff{eWarmUpCompleted}&: \rg{\allTL}{\lastA\msg{eWarmUpCompleted}{\top}}{\finalA\msg{eCoffeeMakerReady}{\top}}
  \\\eff{eNoWaterError}&: \rg{\allTL}{\lastA\msg{eNoWaterError}{\top}}{\finalA\msg{eCoffeeMakerError}{\I{er} = \Code{NoWaterError}}}
  \\\eff{eNoBeansError}&: \rg{\allTL}{\lastA\msg{eNoBeansError}{\top}}{\finalA\msg{eCoffeeMakerError}{\I{er} = \Code{NoBeansError}}}
\end{align*}}\noindent
When the coffee machine processes the \(\eff{eGrindBeansReq}\) and \(\eff{eStartEspressoReq}\) messages, it non-deterministically sends \(\eff{eNoWaterError}\) and \(\eff{eNoBeansError}\) to the panel. The panel then directly reports these errors to the controller via \(\eff{eCoffeeMakerError}\).

\paragraph{\textsc{BankServer}} The controller simulate an user that interacts with a bank to withdraw money from their accounts, where the balance is stored in another database component.
{\small\begin{align*}
  &\eff{eInitAccount}: \Code{ac{:}aid} \sarr \Code{ba{:}int} \sarr 
  \\&\quad \rg{\neg\finalA\msg{eInitAccount}{\I{accountId} = \Code{ac}}}{\lastA\msg{eInitAccount}{\I{accountId} = \Code{ac} \land \I{balance} = \Code{ba}}}{\allTL}
  \\&\eff{eWithDrawReq}: \Code{id{:}rid} \sarr \Code{ac{:}aid} \sarr \Code{am{:}int} \sarr 
  \\&\quad \effLR{\allTL}\effLR{\lastA\msg{eWithDrawReq}{\I{rId} = \Code{id}\land\I{accountId} = \Code{ac} \land \I{amount} = \Code{am}}}
  \\&\quad \effLR{\finalA\msg{eReadQuery}{\I{rId} = \Code{id}\land\I{accountId} = \Code{ac} \land \I{amount} = \Code{am}}}
  \\&\eff{eWithDrawResp}: \Code{id{:}rid} \sarr \Code{ac{:}aid} \sarr \Code{ba{:}int} \sarr \Code{st{:}status} \sarr 
  \\&\quad \rg{\allTL}{\lastA\msg{eWithDrawResp}{\I{rId} = \Code{id}\land\I{accountId} = \Code{ac} \land \I{balance} = \Code{ba} \land \I{status} = \Code{st}}}{\allTL}
  \\&\eff{eUpdateQuery}: \Code{ac{:}aid} \sarr \Code{ba{:}int} \sarr
  \\&\quad \rg{\allTL}{\lastA\msg{eUpdateQuery}{\I{accountId} = \Code{ac} \land \I{balance} = \Code{ba}}}{\allTL}
  \\&\eff{eReadQuery}: \Code{ba{:}int} \garr \Code{id{:}rid} \sarr \Code{ac{:}aid} \sarr \Code{am{:}int} \sarr 
  \\&\quad \effLR{\finalA (\msg{eInitAccount}{\I{accountId} = \Code{ac} \land \I{balance} = \Code{ba}} \land \neg\finalA \msg{eUpdateQuery}{\I{accountId} = \Code{ac}})}
  \\&\quad \effLR{\lastA\msg{eReadQuery}{\I{rId} = \Code{id}\land\I{accountId} = \Code{ac} \land \I{amount} = \Code{am}}}
  \\&\quad \effLR{\finalA\msg{eReadQueryResp}{\I{rId} = \Code{id}\land\I{accountId} = \Code{ac} \land \I{amount} = \Code{am}}\land \I{balance} = \Code{ba}} \interty
  \\&\quad \effLR{\finalA (\msg{eUpdateQuery}{\I{accountId} = \Code{ac} \land \I{balance} = \Code{ba}} \land \neg\finalA \msg{eUpdateQuery}{\I{accountId} = \Code{ac}})}
  \\&\quad \effLR{\lastA\msg{eReadQuery}{\I{rId} = \Code{id}\land\I{accountId} = \Code{ac} \land \I{amount} = \Code{am}}}
  \\&\quad \effLR{\finalA\msg{eReadQueryResp}{\I{rId} = \Code{id}\land\I{accountId} = \Code{ac} \land \I{amount} = \Code{am}}\land \I{balance} = \Code{ba}}
  \\&\eff{eReadQueryResp}: \Code{id{:}rid} \sarr \Code{ac{:}aid} \sarr \Code{am{:}int} \sarr \Code{ba{:}int} \sarr 
  \\&\quad \effLR{\allTL}\effLR{\lastA\msg{eReadQueryResp}{\I{rId} = \Code{id}\land\I{accountId} = \Code{ac} \land \I{amount} = \Code{am} \land \I{balance} = \Code{ba} \land \neg \Code{ba \geq am}}}
  \\&\quad \effLR{\finalA\msg{eWithDrawResp}{\I{rId} = \Code{id}\land\I{accountId} = \Code{ac} \land \neg\I{status}}} \interty
  \\&\quad \effLR{\allTL}\effLR{\lastA\msg{eReadQueryResp}{\I{rId} = \Code{id}\land\I{accountId} = \Code{ac} \land \I{amount} = \Code{am} \land \I{balance} = \Code{ba}} \land \Code{ba \geq am}}  
  \\&\quad \effLR{\finalA\msg{eWithDrawResp}{\I{rId} = \Code{id}\land\I{accountId} = \Code{ac} \land \I{status}} \land \finalA\msg{eUpdateQuery}{\I{accountId} = \Code{ac} \land \I{balance} = \Code{ba - am}}}
\end{align*}}\noindent
